\clearpage\begingroup\NoSourceHead\SourceLeaf{513b}
\renewcommand{\thefootnote}{\arabic{footnote})}\setcounter{footnote}{5}
\ArithmeticHead{56}{L}{BOLZA, GAUSS UND DIE VARIATIONSRECHNUNG.}
\begin{center}\textls[100]{1) Der Transversalensatz.}\end{center}
\noindent\hspace*{\parindent}a) Der \textsc{Gauss}sche Satz für geodätische Linien.
\par\noindent\hspace*{\parindent}6. In erster Linie haben wir hier den \textsc{Gauss}schen Satz über geodätische
Polarkoordinaten (Art. 15) anzuführen: »Zieht man auf einer Fläche von demselben
Ausgangspunkt aus nach allen Richtungen geodätische Linien von
gleicher Länge, so schneidet die Verbindungskurve ihrer Endpunkte die sämtlichen
geodätischen Linien senkrecht«, und den analogen Satz über geodätische
Parallelkoordinaten (Art. 16): »Zieht man von den Punkten einer auf einer
Fläche gegebenen Kurve unter rechtem Winkel und nach derselben Seite hin
geodätische Linien von gleicher Länge, so schneidet die Verbindungskurve
ihrer Endpunkte die sämtlichen geodätischen Linien senkrecht«. Wie schon
\textsc{Gauss} bemerkt hat, kann man den ersten Satz als Spezialfall des zweiten
auffassen, indem man die gegebene Kurve auf einen Punkt (unendlich kleinen
Kreis) zusammenschrumpfen lässt.

\textsc{Gauss}’ Beweis\footnote{Dieser Beweis findet sich bereits in dem ersten Entwurf zu den \textit{Disquisitiones} vom Jahre 1825,
\textit{Neue allgemeine Untersuchungen über die krummen Flächen}, Art. 17, \textsc{Gauss}’ Werke, Bd. VIII, S. 437.} der beiden Sätze kommt darauf hinaus, dass er auf der
Fläche geodätische Polar-, bezw. Parallelkoordinaten, $r,\varphi$ einführt und dann
zeigt, dass für ein solches System von krummlinigen Koordinaten die (sonst
von \textsc{Gauss} mit $F$ bezeichnete) Grösse
\begin{equation*}\tag{9}
S(r,\varphi)=\frac{\partial x}{\partial r}\frac{\partial x}{\partial\varphi}
+\frac{\partial y}{\partial r}\frac{\partial y}{\partial\varphi}
+\frac{\partial z}{\partial r}\frac{\partial z}{\partial\varphi}
\end{equation*}
verschwindet. Denn die partielle Ableitung
\[
\begin{aligned}
\frac{\partial S}{\partial r}={}&\frac{\partial^2x}{\partial r^2}\frac{\partial x}{\partial\varphi}
+\frac{\partial^2y}{\partial r^2}\frac{\partial y}{\partial\varphi}
+\frac{\partial^2z}{\partial r^2}\frac{\partial z}{\partial\varphi}\\
&+\frac12\frac{\partial}{\partial\varphi}\left(\left(\frac{\partial x}{\partial r}\right)^2+\left(\frac{\partial y}{\partial r}\right)^2+\left(\frac{\partial z}{\partial r}\right)^2\right)
\end{aligned}
\]
verschwindet, weil nach der geometrischen Bedeutung von $r$ als Bogenlänge
einer geodätischen Linie einerseits
\clearpage\SourceLeaf{514}
\ArithmeticHead{57}{R}{DER TRANSVERSALENSATZ.}
\begin{equation*}\tag{10}
\left(\frac{\partial x}{\partial r}\right)^2+\left(\frac{\partial y}{\partial r}\right)^2+\left(\frac{\partial z}{\partial r}\right)^2=1
\end{equation*}
und andererseits auf Grund der Fundamentaleigenschaft (2) der geodätischen
Linien auch die erste Summe verschwindet. Daher ist
\begin{equation*}\tag{11}
S(r,\varphi)=S(0,\varphi),
\end{equation*}
und die letztere Grösse ist null, wie sich für den zweiten Satz aus der Konstruktion,
für den ersten aus dem oben angedeuteten Grenzübergang ergibt.
\par\noindent\hspace*{\parindent}b) Der \textsc{Kneser}sche Transversalensatz.
\par\noindent\hspace*{\parindent}7. Bei dem Versuch, diese Sätze auf das allgemeine Integral (8) auszudehnen,
kommt nun alles darauf an, was man als Verallgemeinerung des
senkrechten Schnitts einer geodätischen Linie durch eine schneidende Kurve
aufzufassen hat. \textsc{Kneser}\footnote{A. \textsc{Kneser}, \textit{Lehrbuch der Variationsrechnung}, Braunschweig 1900, § 15.} hat als solche den \textls[100]{transversalen} Schnitt der
Extremale durch eine schneidende Kurve erkannt und dementsprechend als
Verallgemeinerung des zweiten \textsc{Gauss}schen Satzes den \textls[100]{Transversalensatz}
erhalten: »Schneidet man auf den verschiedenen Extremalen einer Extremalenschar
von ihren Schnittpunkten mit einer Transversale $\mathfrak C_0$ aus in positiver
Richtung Bogen ab, die für das Integral $J$ denselben konstanten Wert liefern,
so ist die Verbindungskurve der Endpunkte dieser Bogen wieder eine Transversale
der Schar«, und einen entsprechenden Satz für den speziellen Fall,
wo die Ausgangskurve $\mathfrak C_0$ auf einen Punkt zusammenschrumpft. \textsc{Kneser} geht
beim Beweis dieses Satzes von der Betrachtung eines Extremalenfeldes aus
und leitet den Satz aus den weiter unten zu besprechenden \textsc{Hamilton}schen
Formeln (24) für die partiellen Ableitungen des Feldintegrals ab. Der Satz
ist jedoch von der Voraussetzung, dass die betrachtete Extremalenschar ein
Feld bildet, unabhängig und lässt sich ohne Benutzung der Formeln (24) durch
eine unmittelbare Verallgemeinerung der \textsc{Gauss}schen Schlussweise beweisen.

Es sei in der Tat eine Extremalenschar mit dem Parameter $v$ gegeben und
eine Kurve $\mathfrak C_0$, die jede Extremale der Schar in einem Punkt schneidet. Wenn
dann die Funktion $F$ entlang den Extremalen der Schar positiv ist, so kann
man auf der Extremale $v$ von ihrem Schnittpunkt $Q$ mit der Kurve $\mathfrak C_0$ aus
nach der positiven Seite hin einen Bogen $QP$ abgrenzen, für den das Integral
$J$ einen vorgeschriebenen positiven Wert $u$ annimmt. Die Koordinaten $x,y$
\par\smallskip\noindent{\small X 2 Abh. 5.}\hfill{\small 8}\par
\clearpage\SourceLeaf{515}
\ArithmeticHead{58}{L}{BOLZA, GAUSS UND DIE VARIATIONSRECHNUNG.}
\noindent des Punktes $P$ sind dann eindeutige Funktionen von $u$ und $v$
\begin{equation*}\tag{12}
x=X(u,v),\quad y=Y(u,v),
\end{equation*}
und nach der Bedeutung von $u$ ist, der Gleichung (10) entsprechend,
\begin{equation*}\tag{13}
F(X,Y,X_u,Y_u)=1.
\end{equation*}

Jetzt bilde man den Ausdruck
\begin{equation*}\tag{14}
S(u,v)=F_{x'}(X,Y,X_u,Y_u)X_v+F_{y'}(X,Y,X_u,Y_u)Y_v,
\end{equation*}
dessen Verschwinden ausdrückt, dass die durch den Punkt $P$ gehende Kurve
$u=\text{const.}$ im Punkte $P$ die Extremale $QP$ transversal schneidet. Dann ist
\begin{equation*}\tag{15}
\frac{\partial S}{\partial u}=\frac{\partial F}{\partial v}
-\left[X_v\left(F_x-\frac{\partial}{\partial u}F_{x'}\right)+Y_v\left(F_y-\frac{\partial}{\partial u}F_{y'}\right)\right].
\end{equation*}
Da die Kurve $QP$ Extremale ist, ist dann
\begin{equation*}\tag{16}
\frac{\partial S}{\partial u}=\frac{\partial F}{\partial v},
\end{equation*}
also wegen (13): $\frac{\partial S}{\partial u}=0$ und daher
\begin{equation*}\tag{17}
S(u,v)=S(0,v).
\end{equation*}

Wenn nun insbesondere die Kurve $\mathfrak C_0$ eine Transversale der Extremalenschar
ist, so ist $S(0,v)=0$ und daher allgemein $S(u,v)=0$, womit der Satz
bewiesen ist.
\par\clearpage\endgroup
